\documentclass[11pt]{article}
\usepackage[margin=1in]{geometry}
\usepackage{booktabs}
\usepackage{float}
\usepackage{amsmath}
\usepackage{url}
\usepackage[hidelinks]{hyperref}
\setlength{\emergencystretch}{2em}

\title{DBLP--Scholar Evaluation: Results and Protocol Deviations}
\author{}
\date{2026-09-08}

\begin{document}
\maketitle

With the method configuration frozen before evaluation on DBLP--Scholar,
pair-adaptive selection achieved higher end-to-end F1 than showing all attributes
in both observed execution windows. In the two repetitions per arm on 7 September,
mean F1 was 0.5963 versus 0.5534, with an 18.47\% reduction in pipeline tokens.
The gain came from fewer false positives, accompanied by lower recall.
Blocking recovered only 47.50\% of the gold matches, which limits the entire pipeline.

The planned three-draw same-day experiment was not completed as registered.
One draw per arm completed on 4 September; two per arm completed on 7 September
after changes to recovery and retry handling. The windows are reported separately.
This report supplements the submitted Week 6 report and does not replace its results.

\section{Protocol and Configuration}
DBLP--Scholar contains 2,616 Table-A records, 64,263 Table-B records, 5,347
gold matches and four attributes: title, authors, venue and year. Method settings
were copied from DBLP--ACM before the first evaluation, as recorded in
\path{docs/heldout_preregistration_2026-09-04.md}. The adaptive arm profiles
candidate pairs without using gold labels, fits a transfer model and selects a
subset for each pair. The full arm shows all four attributes. Both arms receive
the same candidate set and use the same matcher prompt template.

\begin{table}[H]
\centering\small
\begin{tabular}{p{0.24\textwidth}p{0.68\textwidth}}
\toprule
Stage & Frozen settings \\
\midrule
Embedding & \texttt{BAAI/bge-large-en-v1.5}; per-column encoding, concatenation and normalisation \\
Blocking & 15 LSH tables, 8 planes, one-bit probing, top-$k=5$ per Table-A record \\
Profiling & 50 pairs; \texttt{diversity} sampler; \texttt{decisive} elicitation; 4 workers \\
Transfer & Ridge ($\alpha=1$), cosine 5-neighbour retrieval, weighted average 0.5/0.5 \\
Selection & \texttt{gap}; minimum 1, maximum 5 (schema width 4); required role \texttt{identity}; mask \texttt{drop}; canonical field order \\
Matcher & \texttt{deepseek/deepseek-v4-flash-0731}; temperature 0; reasoning disabled; output cap 256; 16 workers; unpinned provider routing \\
Seeds & Pipeline 42; blocking 42 \\
\bottomrule
\end{tabular}
\caption{Settings used for the completed Scholar comparison. The current blocking
implementation always probes one-bit neighbours, matching the frozen
\texttt{num\_flips=1}; other values of that argument are not implemented.}
\end{table}

\textbf{Operational deviations.} Attempts on 5 September did not produce valid
results: each draw-2 attempt ended with one empty matcher response, and adaptive
draw 3 stopped during blocking. On 7 September, a separate pair-indexed checkpoint
was added to each remaining draw while response caching remained disabled.
The empty-response limit increased from three to six attempts. Failed-attempt
token usage is now retained. The method settings above were unchanged.
The dated deviations log classifies these as exploratory repetitions under an
amended operational protocol, with frozen method settings.

The four new runs completed between 15:17 and 17:32 Brisbane time on 7 September,
in the order adaptive 2, full 2, adaptive 3, full 3. None required a restart.
All 52,320 new pair answers were saved, with zero cache or checkpoint hits during
execution. One empty response in full draw 3, served by DeepInfra with
\texttt{finish\_reason=stop}, succeeded on retry; its 275 reported tokens are included
in that run's cost. There were 52,321 application-level completion attempts.
SDK-internal retries do not expose complete usage accounting.

\textbf{Validation.} Recalculation from the six CSVs agrees with their run summaries.
All six contain the same 13,080 distinct candidate pairs, no terminal matcher errors
and no response-cache hits. All adaptive profiling error counts are zero. The 52,320
new checkpoint records agree with their CSV answers, attributes, provider, request
hash, raw response and aggregate matching tokens. The regression suite has 153
passing tests, including interrupted-run recovery with mocked calls.

\section{Results by Execution Window}
Candidate recall uses the 2,540 gold matches present in the candidate set as its
denominator. End-to-end recall uses all 5,347 gold matches. Precision is identical
at both levels because the pipeline predicts positives only among candidates.
F1 is computed separately at each level. Thus a high candidate F1 does not imply
high recall over the full benchmark.

\begin{table}[H]
\centering\small
\setlength{\tabcolsep}{4pt}
\begin{tabular}{lrrrrrr}
\toprule
Arm / draw & P & Cand. R & Cand. F1 & E2E R & E2E F1 & Tokens \\
\midrule
Adaptive / 1 & 0.9518 & 0.9559 & 0.9538 & 0.4541 & 0.6148 & 3,045,919 \\
Full / 1     & 0.8896 & 0.9744 & 0.9301 & 0.4629 & 0.6089 & 3,748,500 \\
\bottomrule
\end{tabular}
\caption{4 September: one draw per arm under the original operational protocol.
The end-to-end F1 difference is +0.0059, with 18.74\% fewer pipeline tokens.}
\end{table}

\begin{table}[H]
\centering\small
\setlength{\tabcolsep}{4pt}
\begin{tabular}{lrrrrrr}
\toprule
Arm / draw & P & Cand. R & Cand. F1 & E2E R & E2E F1 & Tokens \\
\midrule
Adaptive / 2 & 0.8754 & 0.9516 & 0.9119 & 0.4520 & 0.5962 & 3,031,488 \\
Full / 2     & 0.7094 & 0.9728 & 0.8205 & 0.4621 & 0.5597 & 3,717,979 \\
Adaptive / 3 & 0.8690 & 0.9559 & 0.9104 & 0.4541 & 0.5965 & 3,031,597 \\
Full / 3     & 0.6834 & 0.9602 & 0.7985 & 0.4561 & 0.5471 & 3,718,319 \\
\midrule
Adaptive mean & 0.8722 & 0.9537 & 0.9111 & 0.4531 & 0.5963 & 3,031,542.5 \\
Full mean     & 0.6964 & 0.9665 & 0.8095 & 0.4591 & 0.5534 & 3,718,149 \\
\bottomrule
\end{tabular}
\caption{7 September: two draws per arm under amended recovery behaviour.
Means are arithmetic means of the two run-level values, not pooled-pair metrics.}
\end{table}

For 7 September, the end-to-end F1 ranges are [0.5962, 0.5965] for adaptive and
[0.5471, 0.5597] for full. The mean difference is +0.0430; the candidate-level
difference is +0.1016. Pipeline tokens fall by 18.47\%. These are two-draw
observations and are not substituted for the registered three-draw same-day test.
No across-date mean is used as the primary comparison.

Pipeline token totals include the fixed profiling workload (15,208 tokens per
adaptive run), even when profiling is replayed from cache. They describe the
logical pipeline workload, rather than an account invoice for each invocation.
The 7 September figures also include reported failed-attempt usage; the older
logging did not retain that overhead.

\section{What Produces the Gain}
\begin{table}[H]
\centering\small
\begin{tabular}{lrrrr}
\toprule
Draw & Adaptive TP & Full TP & Adaptive FP & Full FP \\
\midrule
2 & 2,417 & 2,471 & 344 & 1,012 \\
3 & 2,428 & 2,439 & 366 & 1,130 \\
\bottomrule
\end{tabular}
\caption{Errors on the same candidate sets on 7 September.}
\end{table}

In draw 2, condensation removes 668 false positives while losing 54 true positives.
In draw 3, it removes 764 false positives while losing 11 true positives. The
precision improvement outweighs the recall reduction in the observed F1 scores.
This supports an observed reduction in false positive matching; the experiment
does not identify the model's internal reasoning process.

The selector shows an average of 1.3446 attributes per pair, compared with four
for full. Title appears on all 13,080 pairs, authors on 3,065, year on 1,442 and
venue on none. The two new adaptive runs reuse all 50 profiling responses, and
their selected subsets are identical to draw 1. These repetitions therefore
measure matcher variability conditional on one fixed profile sample and selector;
they do not estimate variability from resampling or re-profiling.

\textbf{Blocking limits coverage.} The common candidate set contains 2,540 of
5,347 gold matches. The remaining 2,807 never reach either matcher arm. Even a
perfect matcher on these candidates could attain only 47.50\% end-to-end recall.
The completed comparison must retain this result. Tuning blocking on Scholar
would be a separate exploratory study, not a revision of this frozen evaluation.

\section{Stability and Prediction Audit}
Between draws 2 and 3, adaptive changes its answer on 565/13,080 pairs (4.32\%),
whereas full changes on 1,754/13,080 (13.41\%). Adaptive is more stable in this
observed pair of repeats. Two runs per arm do not establish a general variance
reduction claim. Provider mixtures also differ between runs; for example,
OpenInference accounts for 26.2\% of full draw 2 and 34.7\% of full draw 3.
Same-day execution reduces the date confound but does not fix the serving mixture
or isolate routing from within-provider sampling. The much different F1 levels
on 4 and 7 September must not be attributed to a change in attribute selection.

\begin{description}
\item[P1: higher F1 and disjoint three-draw ranges.]
The positive direction is supported within both observed windows. The registered
three-draw same-day test was not completed as specified, so P1 is not marked as
fully confirmed. The amended two-draw ranges on 7 September are disjoint.
\item[P2: F1 gain between +0.05 and +0.15.]
The 7 September end-to-end gain is +0.0430, a partial miss. End-to-end F1 is used
here consistently with the DBLP--ACM +0.186 comparison motivating this prediction;
the larger candidate-level gain is not substituted after seeing the outcome.
\item[P3: 15--35\% token reduction.]
The observed 18.47\% reduction on 7 September falls within the interval.
\item[P4: 1.2--2.5 attributes per pair.]
The observed mean of 1.3446 falls within the interval.
\item[P5: title most frequent and present on over half the pairs.]
Numerically satisfied, but required-role enforcement forces the only identity
attribute, title, into the selection. This is not independent evidence that
the learned importance model discovers title's value.
\item[P6: more mask activity than on DBLP--ACM.]
An offline reconstruction of the frozen selector reproduces every observed
subset and finds 120 pairs whose selected subset changes when the mask is applied,
compared with the previously reported five on DBLP--ACM. There are 1,453 candidate
pairs with at least one blank, which is a different quantity. The anticipated
year/venue explanation is not established by this count alone. No optional
Scholar mask-ablation matcher runs were performed.
\end{description}

The supported thesis statement is that the frozen adaptive configuration improves
F1 and reduces token workload in the observed Scholar comparisons, mainly through
fewer false positives, at a recall cost. Report this together with the low blocking
coverage, limited repetitions and unresolved service variation. Claims of universal
improvement, provider causation or complete confirmation of the original protocol
are not supported by these results.

\section{Reproducibility and Next Work}
The submitted Week 6 files remain unchanged. The current handoff is
\path{docs/HANDOFF_step15.md}; the dated protocol deviations are appended to the
pre-registration. The offline validator is
\path{examples/analyze_heldout_runs.py}, with output at
\path{logs/heldout/validated_results.json}. It checks result completeness, errors,
candidate-set equality and agreement between CSV-derived metrics and summaries.
The mask counts and answer-flip counts above were checked separately by offline
selector replay and pair-indexed CSV comparison on 8 September.

Run summaries and CSVs are under \path{logs/DBLP-Scholar/runs/}, in these directories:
\begin{itemize}\small
\item \path{20260904_143805_step15_diversity_decisive_gap_mink1}
\item \path{20260904_151442_full}
\item \path{20260907_155012_step15_diversity_decisive_gap_mink1}
\item \path{20260907_162452_full}
\item \path{20260907_165612_step15_diversity_decisive_gap_mink1}
\item \path{20260907_173241_full}
\end{itemize}

The next planned experiment is the sampler comparison on development datasets,
after explicitly aligning its commands with the headline decisive/gap/minimum-one
configuration. Provider-conditioned work remains a separate open question.
Neither that sweep nor any blocking tuning was run for this report. Cascade work
remains parked under the supervisor's instruction.

\end{document}
